\chapter{Tubes, shadings, and nonconcentration}
\label{chap:geometry}

We work in $\R^3$ unless another dimension is specified. Lebesgue measure
is denoted by $|\cdot|$. The notation $\cE_r(A)$ denotes the least number
of radius-$r$ balls needed to cover $A$, and $N_r(A)$ denotes its
$r$-neighbourhood. Fixed changes in the radius affect covering numbers
by a dimensional constant. We write $A\lesssim B$ when the implicit
constant is independent of the fine scale $\delta$, and $A\lessapprox B$
when a fixed power of $\log(1/\delta)$ is also allowed. Every use of such
a loss takes place before the final small-scale threshold is chosen.

\section{The geometric objects}

\begin{definition}[Tubes; WZ2, Section 2.1]
\label{def:tube}
\lean{Kakeya.DeltaTube,Kakeya.Assouad.WZ2PaperOrdinaryIsEssentiallyDistinct}
\leanok
A $\delta$-tube is the closed $\delta$-neighbourhood of a unit line
segment. Its direction is the unoriented direction of the segment. If
$T$ is a tube and $a>0$, then $aT$ is its dilation by $a$ about its
midpoint. Two tubes $T,T'$ are essentially distinct when
$T\not\subset 2T'$ and $T'\not\subset 2T$.
\end{definition}

\begin{definition}[Shadings; WZ2, Definition 2.1]
\label{def:shading}
\lean{Kakeya.Streamlined.Shading,Kakeya.Streamlined.Shading.IsLambdaDense}
\leanok
\uses{def:tube}
A shading of a finite family $\mathcal S$ is a collection of measurable
sets $Y(S)\subset S$. Write
\[
 E_Y=\bigcup_{S\in\mathcal S}Y(S),\qquad
 \mass(Y)=\sum_{S\in\mathcal S}|Y(S)|,\qquad
 \mathcal S_Y(p)=\{S\in\mathcal S:p\in Y(S)\}.
\]
The shading is $\lambda$-dense if
$\mass(Y)\geq\lambda\sum_{S\in\mathcal S}|S|$.
It is uniformly $\lambda$-dense if $|Y(S)|\geq\lambda|S|$ for every
$S$. The multiplicity at $p$ is $\#\mathcal S_Y(p)$.
\end{definition}

\begin{remark}[Density and multiplicity]
\label{rem:aggregate-density}
The endpoint uses aggregate density, exactly as in WZ2 Definition 2.1.
Uniform density is obtained by discarding tubes with small shading and
then refining. It must not be inserted as an extra hypothesis of the
main theorem. Also, $\mathcal S_Y(p)$ counts shaded incidences: replacing
$p\in Y(S)$ by $p\in S$ changes the multiplicity estimates used below.
The families in Lean are indexed, so multiplicities are retained until
an explicit selection of representatives is made.
\end{remark}

\begin{definition}[Regular shadings; WZ2, Definition 2.2]
\label{def:regular-shading}
\uses{def:shading}
For a measurable set $S\subset B(0,1)\subset\R^n$ of positive measure,
a shading is regular down to scale $\delta$ if, for $x\in Y(S)$ and
$\delta\leq r\leq1$,
\[
 |Y(S)\cap B(x,r)|\geq
 \frac{|Y(S)|}{100^n\log(1/\delta)}
 \frac{|S\cap B(x,r)|}{|S|}.
\]
Here the fine scale is sufficiently small that the denominator is
positive.
\end{definition}

\begin{lemma}[Regularization; WZ2, Lemma 2.3]
\label{lem:regular-shading}
\uses{def:regular-shading}
Given a shading $Y(S)$ as above, one can pass to a regular subshading
$Y'(S)\subset Y(S)$ while retaining at least half its measure.
\end{lemma}
\begin{proof}
At each dyadic scale, discard the cells where the shaded measure is
too small relative to the measure of $S$. The threshold allocates only
a fixed fraction divided by the number of scales to each deletion.
The sum of the discarded measures is at most half the original measure.
A bounded collection of shifted dyadic grids supplies the ball version.
\end{proof}

\section{Wolff conditions}

\begin{definition}[Wolff axioms; WZ2, Definition 2.4]
\label{def:wolff}
\uses{def:tube}
A family of $\delta$-tubes in $\R^n$ satisfies the Wolff axioms with
constant $C$ if every $\rho$-tube contains at most
$C(\rho/\delta)^{n-1}$ members, and every rectangular prism of dimensions
$2\delta\times\cdots\times2\delta\times\rho\times2$ contains at most
$C\rho/\delta$ members. The conditions are imposed for
$\delta\leq\rho\leq1$.
\end{definition}

\begin{definition}[Convex Wolff Axioms; WZ2, Definition 2.5]
\label{def:cwa}
\lean{Kakeya.Assouad.WZ2PaperBodyConvexWolffBound}
\leanok
A finite indexed family $\mathcal S$ satisfies the Convex Wolff Axioms
with error $C$ if every convex set $W\subset\R^3$ obeys
\[
  \#\{S\in\mathcal S:S\subset W\}\leq C|W|\,\#\mathcal S.
\]
We always take $1\leq C<\infty$.
\end{definition}

\begin{remark}[Consequences; WZ2, Remarks 2.6--2.7]
\label{rem:cwa-consequences}
If $\mathcal S$ is nonempty and its members are convex, applying the
axiom with $W=S$ gives
$\#\mathcal S\geq C^{-1}(\inf_{S\in\mathcal S}|S|)^{-1}$.
For equal-radius tubes in three dimensions, this is
$\#\T\gtrsim C^{-1}\delta^{-2}$.
If every member is replaced by a larger set, with its index retained,
the axiom continues to hold: containment of the larger set implies
containment of the original one. Identifying repeated larger sets is
a separate operation and requires a counting argument.
\end{remark}

\begin{definition}[Unit rescaling; WZ2, Definition 2.8]
\label{def:unit-rescaling}
\lean{Kakeya.Assouad.WZ2PaperAssouadUnitRescalingData}
\leanok
Let $U$ be a convex body, and let $U^*$ be its minimum-volume
circumscribed ellipsoid. Choose an affine isomorphism $\phi_U$ carrying
$U^*$ onto the unit ball. The unit rescaling of a family of subsets of
$U$ is the indexed family of their images under $\phi_U$.
\end{definition}

\begin{remark}[The two normalizations; WZ2, Remark 2.9]
\label{rem:john-normalization}
WZ1 uses a closely related normalization adapted to its cubical
shadings; the terminal ellipsoid has fixed, comparable axes rather
than being exactly the unit ball. WZ2 uses the actual circumscribed
John ellipsoid. The public Lean definition follows WZ2. In particular,
the rescaled family consists of affine images of the original sets,
not arbitrarily chosen nearby tubes. Replacing these images by
standard tubes later in the argument includes explicit containment
and volume comparisons.
\end{remark}

\begin{definition}[Uniform partitioning covers; WZ2, Definition 2.10]
\label{def:partition-cover}
\lean{Kakeya.Assouad.WZ2PaperPurePartitioningCover,Kakeya.Assouad.WZ2PaperPureFullFibersAreCUniform}
\leanok
\uses{def:tube}
For families $\mathcal A,\mathcal B$ set
$\mathcal B[A]=\{B\in\mathcal B:B\subset A\}$.
The family $\mathcal A$ covers $\mathcal B$ if every member of
$\mathcal B$ is contained in some $A\in\mathcal A$. It is a
$C$-uniform cover if
\[
 \#\mathcal B[A]\leq C\#\mathcal B[A']\quad(A,A'\in\mathcal A).
\]
It is partitioning if the indexed sets $\mathcal B[2A]$ and
$\mathcal B[2A']$ are disjoint whenever $A\ne A'$.
\end{definition}

\begin{remark}[Geometric fibres]
\label{rem:full-fibres}
The fibre $\mathcal B[A]$ is the complete containment fibre. One cannot
replace it by a selected group of children assigned to $A$ and then
assert the same uniformity property. In the Lean definition the unique
parent map is derived from containment and doubled-fibre disjointness.
It is not an independent choice. This distinction matters after
refining either family.
\end{remark}

\begin{lemma}[Passing CWA to a cover; WZ2, Remark 2.11]
\label{lem:cover-cwa}
\uses{def:cwa,def:partition-cover}
If $\mathcal B$ is nonempty and satisfies CWA with error $C$, and
$\mathcal A$ is a $K$-uniform partitioning cover, then $\mathcal A$
satisfies CWA with error $KC$.
\end{lemma}
\begin{proof}
The nonempty fibres partition $\mathcal B$. Their minimum size is at
least $\#\mathcal B/(K\#\mathcal A)$. Every $A\subset W$ contributes
at least this many distinct members of $\mathcal B$ contained in $W$.
Apply CWA to $\mathcal B$ and cancel $\#\mathcal B$.
\end{proof}

\begin{definition}[CWA at every scale; WZ2, Definition 2.12]
\label{def:all-scale-cwa}
\lean{Kakeya.Assouad.WZ2PaperPureCWAAtNearbyScales}
\leanok
\uses{def:cwa,def:partition-cover,def:unit-rescaling}
A family $\T$ of essentially distinct $\delta$-tubes satisfies CWA at
every scale with error $C$ if, for every requested scale
$\rho_0\in[\delta,1]$, there is a scale
$\rho\in[\rho_0,C\rho_0)$ and a family $\T_\rho$ of $\rho$-tubes such
that $\T_\rho$ is a $C$-uniform partitioning cover of $\T$, and the
unit rescaling of $\T[T_\rho]$ satisfies CWA with error $C$ for each
$T_\rho\in\T_\rho$.
\end{definition}

\begin{remark}[Requested and selected scales]
\label{rem:nearby-scale}
Only the requested scale $\rho_0$ is constrained to be at most $1$.
The selected scale can exceed $1$. This is relevant when the definition
is applied at the top scale: one cannot silently treat its output as
a family of radius-$1$ tubes. Also, the half-open interval in the
definition is empty at $C=1$. In applications $C=\delta^{-\eta}>1$,
with $0<\delta<1$ and $\eta>0$.
\end{remark}

\begin{lemma}[Multiscale inheritance; WZ2, Lemma 2.13]
\label{lem:multiscale-cwa}
\uses{def:all-scale-cwa}
At any requested scale $\rho_0\in[\delta,1]$, a family satisfying
all-scale CWA with error $C$ admits a nearby coarse scale
$\rho\in[\rho_0,C\rho_0]$ for which both the coarse family and the
unit-rescaled fibres satisfy all-scale CWA with error $O(C)$.
\end{lemma}
\begin{proof}
Choose partitioning covers at a separated sequence of scales. The
disjointness condition on doubled fibres makes the relevant containment
classes nested. The resulting tree has comparable class sizes at each
level. Restricting this tree above or below the selected level gives
the two conclusions. When standard tubes are substituted for affine
images, the fixed geometric comparison factors must be included in
the error constant.
\end{proof}

\begin{definition}[Tube Frostman condition; WZ2, Definition 2.14]
\label{def:tube-frostman}
\lean{Kakeya.Assouad.TubeWolffBound}
\uses{def:tube}
A family of essentially distinct tubes satisfies a Frostman condition
of dimension $s$ with error $C$ if
\[
 \#\T[T_\rho]\leq C\rho^s\#\T
 \qquad(\rho\geq\delta)
\]
for every $\rho$-tube $T_\rho$. The dimension-$2$ condition, without a
prescribed total cardinality, is the Tube Wolff condition used in WZ2
Section 7. Its parameter-space implementation is discussed in
Remark~\ref{rem:parameter-frostman}.
\end{definition}

\begin{definition}[Self-similar CWA; WZ2, Definition 2.15]
\label{def:self-similar-cwa}
\uses{def:all-scale-cwa}
Write $\#\T=\delta^{-s}$. The self-similar Convex Wolff Axioms add to
Definition~\ref{def:all-scale-cwa} the estimates
\[
 C^{-1}(\rho/\delta)^s\leq\#\T[T_\rho]
 \leq C(\rho/\delta)^s
\]
for each selected coarse tube. This stronger condition explains the
relation with stickiness, but is not a hypothesis of our endpoint.
\end{definition}

\section{Grain terminology}

\begin{definition}[Plane maps; WZ1, Definition 4.1]
\label{def:plane-map}
\lean{Kakeya.Assouad.WZ1PlaneMapData}
\uses{def:shading}
A plane map assigns a unit normal $V(p)\in S^2$ to each point of the
shaded union. For every shaded incidence $p\in Y(T)$ it satisfies
$|V(p)\cdot\dir(T)|\lesssim\delta$. Thus the directions through $p$
are nearly contained in the plane normal to $V(p)$.
\end{definition}

\begin{definition}[Upper AD sets; WZ1, Definition 4.2]
\label{def:ad-set}
\lean{Kakeya.Assouad.PureWZ2PaperADSet1,Kakeya.Assouad.IsADSet1}
\leanok
A bounded set $A\subset\R^n$ is a $(\delta,\alpha,C)_n$-ADset if,
whenever $\delta\leq\rho\leq r$,
\[
 \cE_\rho(A\cap B(x,r))\leq C(r/\rho)^\alpha.
\]
The paper-facing one-dimensional Lean formulation uses intervals in
place of balls and assumes $0<\alpha\leq1$. The internal bounded
formulation restricts the radii to at most $1$ and fixes a bounded
coordinate interval; the two forms are compared with a fixed change
in the covering constant.
\end{definition}

\begin{remark}[AD does not mean a two-sided measure bound here]
\label{rem:upper-ad}
As observed in WZ1 Remark 4.1, this is an upper covering condition and
is inherited by every subset. No lower Ahlfors--David regularity is
asserted. This is essential because many subsequent steps discard
parts of the shading. The empty projected set also satisfies the
covering inequalities.
\end{remark}
